\section*{Future Work}

\subsection*{Program Analysis for Backend Selection and Overhead Reduction}

The largest overhead driver is the number of write faults fired per unit of
work, which is determined entirely by each kernel's memory access pattern.
The existing \texttt{dirty\_tracking\_hint} procfs entry already lets an
operator select between sequential (sorted-vector) and random (nested-bitmap)
backends before a session, but the selection is manual.

A natural extension is to derive the hint automatically by analysing the compiled
PTX of each kernel.
Store instructions (\texttt{st.*}) carry address operands whose computation
reveals the access pattern: a unit-stride expression where adjacent threads write
adjacent addresses signals sequential access; irregular or data-dependent address
arithmetic signals random access.
This classification feeds directly into the existing hint dispatch, eliminating
operator involvement.

Beyond backend selection, alias analysis on the NVVM IR (LLVM dialect emitted
before device linking) can identify which kernel arguments are derived from
\texttt{cudaMallocManaged} pointers.
At \texttt{start} and \texttt{resume}, only VA blocks reachable through those
arguments need their write permissions revoked.
The current implementation walks the entire owner VA space and calls
\texttt{uvm\_va\_block\_revoke\_prot\_mask} on every mapped block; restricting
this walk to write-reachable blocks reduces the downgrade barrier latency
proportionally to the fraction of the address space that is actually written.
This is most valuable at large working-set sizes (32-64\,GB) where the
downgrade walk is a non-trivial fraction of \texttt{start}/\texttt{resume}
latency.

\subsection*{Migration Policy for the Oversubscribed Regime}

In the oversubscribed regime the GPU's on-device memory is exhausted and UVM
evicts pages to host DRAM using an LRU policy that has no knowledge of write
intensity.
The dirty tracking data already produced by the existing implementation is
sufficient to improve several aspects of this policy.

\paragraph{Write-heat aware eviction}
In cumulative tracking mode, the dump accumulates which pages have been dirtied
across multiple epochs.
A page that appears dirty in every epoch is rewritten every time it enters GPU
memory; evicting it immediately triggers a re-fault.
A page absent from the dirty set for several epochs is effectively read-only
and can be evicted without a write-back.
Adding an epoch-frequency counter per VA block, maintained by user space
from successive dump outputs, and using it to bias UVM's eviction candidate
scoring would keep write-hot pages on-device longer and preferentially evict
clean pages, reducing re-fault frequency under memory pressure.

\paragraph{Cutover timing for iterative migration}
In iterative pre-copy live migration, the migration controller issues
\texttt{query\_cutover} to snapshot the dirty set and then transfers those
pages to the destination.
Convergence requires that the re-dirty rate during transfer is lower than the
transfer bandwidth.
The nanosecond-resolution timestamps in the dump make the write rate directly
observable: the density of timestamps over the epoch's duration gives an estimate
of faults per second.
A migration controller that reads the insert count from \texttt{dirty\_ds\_stats}
in real time can defer \texttt{query\_cutover} until the observed write rate
drops below available transfer bandwidth, improving round-trip convergence
without any additional kernel instrumentation.

\paragraph{Dirty ratio as a scheduling signal}
After each dump, the number of output lines divided by the total managed
allocation size gives the fraction of pages dirtied in that epoch.
When this ratio is low the workload is in a read-intensive phase and migration
rounds can be scheduled more aggressively; when it is high the workload is
producing large dirty sets and extending the inter-round interval allows it to
reach a quieter phase before the next cutover.
This feedback loop is entirely implementable in user space using the existing
procfs interface.
